\section*{Chapitre \ref{ch:b3:structural-biology} --- Biologie structurale des protéines}

\begin{solution}{exo:b3:structural-biology:1}
Primaire~: la suite des $141$ résidus d'une chaîne $\alpha$.
Secondaire~: ses huit hélices $\alpha$. Tertiaire~: le \omterm{def:b3:structural-biology:levels}{repliement}
globine d'une chaîne enroulée autour de son hème. Quaternaire~: le
tétramère $\alpha_{2}\beta_{2}$ et la façon dont ses dimères tournent
l'un contre l'autre.
\end{solution}

\begin{solution}{exo:b3:structural-biology:2}
$36\times \qty{0.15}{nm} = \qty{5.4}{nm}$~; $36/3.6 = 10$ tours~;
$36 - 4 = 32$ liaisons hydrogène.
\end{solution}

\begin{solution}{exo:b3:structural-biology:3}
Que la structure native, y compris l'appariement correct de quatre
ponts disulfure parmi $105$ possibilités, se retrouve spontanément à
partir de la chaîne dépliée~: la séquence contient toute l'information
nécessaire, et l'état natif est le minimum thermodynamique. Le faire
dans un tube, avec de la protéine pure, excluait toute matrice, toute
enzyme et toute machinerie cellulaire comme source de cette
information.
\end{solution}

\begin{solution}{exo:b3:structural-biology:4}
RMN~: petites protéines, sous \qty{40}{kDa} environ, en solution.
Cristallographie~: toute taille qui cristallise, du peptide au
ribosome. \omterm{def:b3:structural-biology:methods}{Cryomicroscopie électronique}~: grands complexes, au-dessus
de \qty{50}{kDa} environ, sans cristaux. La RMN rend compte du
mouvement directement (et la cryo-EM sait trier les molécules en
classes conformationnelles)~; une structure cristalline est une moyenne
statique.
\end{solution}

\begin{solution}{exo:b3:structural-biology:5}
$2^{150} \approx 1.4\times 10^{45}$ conformations, $1.4\times 10^{33}$
s $\approx 5\times 10^{25}$ ans. $3^{20} \approx 3.5\times 10^{9}$,
$3.5$ ms~: le peptide pourrait explorer exhaustivement en une seconde~;
la protéine ne le pourrait pas dans l'âge de l'univers.
\end{solution}

\begin{solution}{exo:b3:structural-biology:6}
$Y = [\alpha(1+\alpha) + Lc\alpha(1+c\alpha)]/[(1+\alpha)^{2} +
L(1+c\alpha)^{2}]$. $\alpha = 1$~: $3.01/106 = 0.028$ (site unique
$0.50$). $\alpha = 10$~: $121/242 = 0.50$ ($0.91$). $\alpha = 100$~:
$\num{10300}/\num{10601} = 0.97$ ($0.99$). La courbe MWC reste loin
derrière le site unique à faible ligand, puis monte abruptement --- de
$0.03$ à $0.97$ sur deux décades là où le site unique passe de $0.5$ à
$0.99$~: cette raideur est la \omterm{def:b3:structural-biology:allostery}{coopérativité}.
\end{solution}

\begin{solution}{exo:b3:structural-biology:7}
Le bisphosphoglycérate ne se fixe qu'à l'\omterm{def:b3:structural-biology:allostery}{état T} (dans la cavité
centrale), donc il stabilise T~: dans le modèle il élève $L$, le
rapport $T/R$ de la protéine sans ligand, et décale donc la courbe vers
la droite --- affinité plus faible, davantage de dioxygène libéré aux
pressions tissulaires. Les hématies conservées perdent ce composé, $L$
baisse, la courbe se décale vers la gauche, et le sang transfusé garde
son dioxygène au lieu de le céder dans les tissus, jusqu'à ce que les
cellules le resynthétisent en une journée.
\end{solution}

\begin{solution}{exo:b3:structural-biology:8}
$d = \lambda/(2\sin\theta) = 0.1/(2\times 0.259) = \qty{0.19}{nm}$~: le
squelette et chaque chaîne latérale sont placés, les atomes individuels
résolus. Pour \qty{0.12}{nm}~: $\sin\theta = 0.1/0.24 = 0.417$, $\theta =
\qty{25}{\degree}$.
\end{solution}

\begin{solution}{exo:b3:structural-biology:9}
Une lysine chargée enfouie parmi des chaînes latérales apolaires coûte
des dizaines de kilojoules par mole (sa charge n'est plus solvatée,
l'enfouissement hydrophobe de la leucine est perdu), plus que tout le
$\Delta G$ du \omterm{def:b3:structural-biology:levels}{repliement}~: la protéine se déplie ou le cœur se
réarrange. En surface, la lysine est hydratée comme elle le serait à
l'état déplié et rien n'est perdu.
$\qty{-30}{kJ/mol} \approx -12RT$ signifie $K = e^{12} \approx 1.6\times
10^{5}$~: une molécule sur \num{160000} est dépliée à tout instant, et
toute la marge vaut deux ou trois liaisons hydrogène sur des centaines
--- de quoi une seule mutation.
\end{solution}

\begin{solution}{exo:b3:structural-biology:10}
Les espèces $R$ totalisent $[R_{0}](1+\alpha)^{n}$ et les espèces $T$
$[R_{0}]L(1+c\alpha)^{n}$, donc $\bar R = (1+\alpha)^{n}/[(1+\alpha)^{n}
+ L(1+c\alpha)^{n}]$. Pour $c \to 0$ le terme $T$ vaut $L$, et $\bar R =
1/2$ lorsque $(1+\alpha)^{n} = L$, c'est-à-dire
$\alpha = L^{1/n} - 1$. Pour l'hémoglobine
$(3\times 10^{5})^{1/4} = 23.4$, donc $\alpha \approx 22$~:
\qty{22}{mmHg}, proche du $P_{50}$ de \qty{26}{mmHg} --- la bascule
conformationnelle et la demi-saturation coïncident presque.
\end{solution}

\begin{solution}{exo:b3:structural-biology:11}
Elle a rencontré de la résistance parce que tout agent infectieux connu
se répliquait par un acide nucléique, et qu'une protéine ne sait pas
recopier sa séquence. Un acide nucléique dont on montrerait qu'il est
nécessaire, ou une protéine purifiée dont on montrerait qu'elle ne
transmet pas, la réfuteraient. À son appui~: l'infectiosité survit aux
nucléases, à l'ultraviolet et aux rayonnements ionisants à des doses
qui détruisent tout génome, mais elle est détruite par les dénaturants
de protéines et par les protéases~; une PrP recombinante repliée in
vitro en fibrilles transmet la maladie à des animaux, avec les mêmes
propriétés de souche au fil des passages. Un germe est nécessaire parce
que la conversion d'un monomère normal est prohibitivement lente à elle
seule --- le noyau de la structure croisée-$\beta$ doit d'abord se
former, ce qui est l'étape limitante --- alors qu'une extrémité de
fibrille déjà formée convertit rapidement les monomères~: la forme se
recopie par croissance sur matrice.
\end{solution}

\begin{solution}{exo:b3:structural-biology:12}
Beaucoup de séquences se replient en une même structure~: ce qu'un
\omterm{def:b3:structural-biology:levels}{repliement} exige, c'est un patron de positions hydrophobes et polaires,
quelques glycines et prolines aux virages, et les résidus qui font le
travail~; tout le reste en surface peut changer sans pénalité, de sorte
que les mutations s'y accumulent tandis que la sélection préserve le
\omterm{def:b3:structural-biology:levels}{repliement} et la fonction. Les résidus conservés sont donc ceux du cœur
(comme patron plus que comme identités), les résidus catalytiques et de
fixation, et les glycines, prolines et cystéines structurales. Les
méthodes de profil pondèrent exactement ces colonnes et reconnaissent
donc des membres de la famille à \qty{15}{\%} d'identité~; la
conception de protéines suit la même logique à l'envers, en choisissant
un patron hydrophobe pour un \omterm{def:b3:structural-biology:levels}{repliement} cible et en laissant la surface
libre.
\end{solution}

\begin{solution}{pb:b3:structural-biology:1}
\textbf{1.} $0.75\times 150 \approx 112$ résidus en hélice~; $112\times
\qty{0.15}{nm} = \qty{17}{nm}$~; $112/3.6 \approx 31$ tours.
\textbf{2.} $112 - 8\times 4 = 80$ liaisons hydrogène.
\textbf{3.} Masse $150\times 110 = \qty{16500}{Da} = 2.7\times 10^{-20}$
g~; volume $2.7\times 10^{-20}\times 0.73 = 2.0\times 10^{-20}$ cm$^{3}$
$= \qty{20}{nm^3}$~; rayon $(3V/4\pi)^{1/3} = \qty{1.7}{nm}$.
\textbf{4.} Étendue $150\times 0.36 = \qty{54}{nm}$ contre un diamètre
de \qty{3.4}{nm}~: un facteur $16$.
\textbf{5.} Les résidus hélicoïdaux font \qty{75}{\%} de la longueur
étendue, $112\times 0.36 = \qty{40}{nm}$~; l'hélice les comprime à
\qty{17}{nm}, d'un facteur $0.36/0.15 = 2.4$.
\textbf{6.} $40\times 4 = \qty{160}{kJ/mol}$ gagnés, $150\times 0.9 =
\qty{135}{kJ/mol}$ perdus~: net $\Delta G \approx \qty{-25}{kJ/mol}$,
dans la fourchette marginale.
\textbf{7.} $3^{150} \approx 4\times 10^{71}$~; $4\times 10^{59}$ s
$\approx 10^{52}$ ans.
\textbf{8.} $10^{-2}\times 10^{12} = 10^{10}$ conformations, soit une
fraction $3\times 10^{-62}$ du compte.
\textbf{9.} Les hélices en \qty{1}{\micro s} (en parallèle)~;
$8! = \num{40320}$ empilements à $10^{6}$ par seconde prennent
\qty{40}{ms}~: au total environ \qty{40}{ms}, l'ordre de grandeur
observé. Oui~: résoudre les parties indépendamment puis le tout est un
entonnoir --- l'espace de recherche est factorisé, non épuisé.
\textbf{10.} Nombre espéré de cycles $1/0.3 = 3.3$, soit environ
\qty{33}{s}~; après six cycles, $0.7^{6} = 0.12$ reste déplié.
\textbf{11.} $\Delta G = -RT\ln K$ avec $K = 10^{4} \to 10^{2}$~: de
$-23.7$ à \qty{-11.9}{kJ/mol}, une perte de \qty{12}{kJ/mol} de
stabilité ($RT = \qty{2.58}{kJ/mol}$, $\ln 100 = 4.6$).
\textbf{12.} L'agrégation commence par un noyau, dont la vitesse de
formation croît comme une puissance élevée de la concentration de
l'espèce dépliée apte à s'agréger~; cent fois plus de cette espèce
élève de plusieurs ordres de grandeur la chance qu'un germe se forme au
cours d'une vie, et une fois le germe présent la croissance est rapide
et se recopie elle-même.
\textbf{13.} $\alpha = 100$~: $Y = 1.054\times 10^{8}/1.089\times 10^{8} =
0.97$. $\alpha = 20$~: numérateur $185\,220 + 103\,680 = 288\,900$,
dénominateur $194\,481 + 622\,080 = 816\,561$~: $Y = 0.35$.
\textbf{14.} $\bar R = 1.041\times 10^{8}/1.089\times 10^{8} = 0.96$ au
poumon~; $194\,481/816\,561 = 0.24$ dans le muscle.
\textbf{15.} $0.97 - 0.35 = 0.62$ de la capacité cédée (les deux tiers
de la charge). Site unique~: $100/126 - 20/46 = 0.79 - 0.43 = 0.36$.
\textbf{16.} Capacité $150\times 1.34 = \qty{200}{mL/L}$~; au muscle
$0.62\times 200 \approx \qty{125}{mL/L}$. Au repos $(0.97 - 0.78)\times
200 = \qty{38}{mL/L}$~; $250/38 \approx \qty{6.5}{L/min}$ de débit
sanguin --- l'ordre d'un débit cardiaque de repos.
\textbf{17.} $L = 3\times 10^{6}$, $\alpha = 40$~: numérateur $2.76\times
10^{6} + 3.29\times 10^{6} = 6.05\times 10^{6}$, dénominateur $2.83\times
10^{6} + 11.5\times 10^{6} = 14.4\times 10^{6}$~: $Y = 0.42$ au lieu de
$0.78$ --- bien plus de dioxygène cédé au repos, adaptation du sang à
l'altitude et du muscle au travail.
\textbf{18.} Fœtale, $L = 3\times 10^{4}$, $\alpha = 30$~: numérateur
$893\,730 + 19\,770 = 913\,500$, dénominateur $923\,520 + 85\,680 =
1\,009\,200$~: $Y = 0.91$. Maternelle à la même pression~: numérateur
$893\,730 + 197\,730 = 1\,091\,460$, dénominateur $923\,520 + 856\,830
= 1\,780\,350$~: $Y = 0.61$. À pression égale le sang fœtal retient
davantage de dioxygène, qui passe donc de la mère au fœtus à travers le
placenta.
\textbf{19.} $(10^{5}\,\text{nm}/6\,\text{nm})^{3} = 4.6\times 10^{12}$
mailles~; $1.9\times 10^{13}$ molécules.
\textbf{20.} $d = 0.1/(2\sin 14.5^{\circ}) = \qty{0.20}{nm}$~: oui,
chaque chaîne latérale et chaque atome sont placés.
\textbf{21.} $\sin\theta = 0.1/0.3$, $\theta = \qty{19.5}{\degree}$.
Dans un cristal désordonné les molécules ne sont pas exactement en
registre, les ondes diffusées aux grands angles --- celles qui codent
le détail fin --- ne s'ajoutent plus en phase, et les taches
s'effacent aux grands angles~; seules survivent les taches de basse
\omterm{def:b3:structural-biology:methods}{résolution}.
\textbf{22.} Diviser $d$ par deux demande deux fois le signal, donc
quatre fois les particules~: environ \num{40000} (davantage en
pratique, d'autres erreurs entrant en jeu).
\textbf{23.} Les protéines membranaires, qui cristallisent rarement à
partir de solutions de détergent et qu'on peut imager dans un
nanodisque lipidique~; les grands complexes souples --- ribosomes,
spliceosomes, canaux ioniques avec leurs régulateurs --- dont
l'hétérogénéité empêche la cristallisation mais dont les particules se
trient en classes conformationnelles.
\textbf{24.} L'hème et le dioxygène fixés et leur géométrie exacte, les
eaux et les ions ordonnés, la seconde conformation (T contre R) et la
preuve que le \omterm{def:b3:structural-biology:levels}{repliement} est juste --- une prédiction est un modèle
unique, sans ligand et sans contrôle expérimental.
\textbf{25.} Un déchargement de $0.62$ de la capacité (site unique
$0.36$)~; le temps de Levinthal d'environ $10^{52}$ ans~; une carte à
\qty{0.20}{nm}.
\end{solution}
